\chapter{SSA Construction}
\label{ch:pebble-ssa}

The source name \code{total} can denote a succession of values during a loop.
An optimizer needs to distinguish those values without repeatedly searching for assignments that might reach a use.
\term{Static single assignment}, or SSA, gives each static definition its own name and makes path-dependent selection explicit.

\section{Direct construction for expressions}
\label{sec:pebble-direct-ssa}
The expression front end constructs SSA directly.
Each arithmetic operation creates a fresh value ID and returns it to its parent.
All references to \code{x} use the single input ID.
There are no source assignments, so neither renaming stacks nor dominance-frontier placement are necessary for this client.
This is a useful first implementation: do not build a general conversion pass before a client needs it.

Branches and loops in the test fixtures are also built directly, with explicit phis.
The executable library verifies these graphs and executes them, but does not yet convert an arbitrary mutable-variable CFG to SSA.
The construction algorithm below specifies that extension and shows what additional state it requires.

\section{Merge a loop-carried value}
\label{sec:pebble-loop-ssa}
The summation fixture computes $\sum_{i=0}^{x-1}i$ modulo 256.
Entry creates input $v_1$, zero $v_2$, and one $v_3$.
An additional constant at $v_4$ is deliberately unused in this fixture and is removed by optimization.
Header values $v_5$ and $v_6$ are the current index and sum.
Their incoming values are
\[
v_5=\phifn{0:v_2,\ 2:v_8},\qquad
v_6=\phifn{0:v_2,\ 2:v_9}.
\]
The header compares $v_5<v_1$ at $v_7$.
The body computes $v_8=v_5+v_3$ and $v_9=v_6+v_5$, then jumps back.
Exit returns $v_6$.

For $x=3$, header states $(i,s)$ are $(0,0)$, $(1,0)$, $(2,1)$, and $(3,3)$.
The last visit exits without executing the body.
The static IDs remain unchanged while the dynamic instances of their values change on each loop traversal.
SSA means one static definition, not one execution of that definition.

\section{The mutable-variable extension}
\label{sec:pebble-ssa-extension}
Algorithm~\ref{alg:pebble-ssa-build} uses iterated dominance frontiers for placement and a dominator-tree walk for renaming.
It assumes definite assignment and a representation whose phi capacity can accommodate every predecessor.
Our two-input representation can implement it only for normalized joins with at most two predecessors; a more general converter must expand the phi representation or transform larger joins first.

\begin{algorithm}[htbp]
\caption{Converting mutable variables to SSA: an extension contract.}
\label{alg:pebble-ssa-build}
\begin{pseudocode}
\Require reachable CFG, definite assignment, dominator tree, frontiers, definition sites
\Ensure renamed graph whose uses satisfy SSA dominance
\ForAll{source variables $v$}
\State seed a placement worklist with the blocks defining $v$
\State insert one phi at each newly reached frontier block
\State treat each inserted phi as a new definition for further placement
\EndFor
\State seed version stacks with entry definitions
\State walk the dominator tree, pushing fresh names at definitions
\State rename ordinary uses from the current stack tops
\State fill each successor phi input from the stack state on its incoming edge
\State pop a block's definitions when returning from its dominator subtree
\State verify the complete renamed graph
\end{pseudocode}
\end{algorithm}

The stack discipline is lexical with respect to the dominator tree, not the source text or block-number order.
If a sibling subtree inherits names pushed by its sibling, its uses may refer to definitions that never execute on that path.
An empty stack at a reachable use is an unfulfilled initialization obligation, not permission to invent a zero constant.

Placement may create dead phis.
Pruned SSA additionally consults liveness, while trivial-phi simplification can remove a merge whose incoming values are identical.
Both are separate refinements of construction.
The current optimizer conservatively retains live phis rather than attempting arbitrary congruence reasoning through cycles.

\section{Use-def information}
\label{sec:pebble-use-def}
PebbleIR's operands are authoritative; it does not maintain reverse use lists.
Replacing all uses therefore scans the table and terminators.
This costs $O(N+B)$, but avoids maintaining two structures during every edit.
If repeated rewrites become expensive, add use records that identify the user, operand position, and, for phis, the associated incoming edge.
Test consistency by rebuilding the reverse relation from operands and comparing it with the maintained lists.

\section{Exercises}
\label{sec:pebble-ten-exercises}
\begin{enumerate}
\item \label{ex:pebble-sum-trace} Trace the summation fixture for inputs 0 and 4, including the final failed test. Give the number of body executions and returned byte.
\item \label{ex:pebble-sibling-stack} In a diamond, the left arm defines \code{a} and the right arm uses the entry definition of \code{a}. Explain the error caused by failing to pop the left arm's renaming stack.
\end{enumerate}

\section{Further reading}
\label{sec:pebble-ten-reading}
The primary SSA construction paper gives the placement and renaming foundations that the extension relies on~\cite{cytron-1991-ssa}.
Read its invariants before implementing a worklist that merely resembles the algorithm.
